\section{Background and Related Work}
\label{sec:background}

\subsection{VeReMi-Extension and Its Sybil Attacks}

VeReMi~\cite{heijden2018veremi} is a simulated dataset of receiver logs with five position-falsification attacks
and no Sybil attack. VeReMi-Extension~\cite{kamel2020veremi}, generated with the F2MD
framework~\cite{kamel2020simulation}, adds a sensor error model and further attacks, among them four Sybil attacks
with the attacker codes A16 to A19. Each log belongs to one receiving vehicle and lists the messages it heard, with the
position, speed, and acceleration claimed by the sender and the sender's pseudonym. In the data used here, each
attack is recorded in two time windows (07--09 and 14--16), and a benign vehicle sends one message per second.

The four attacks differ in what the attacker fabricates and in how it uses pseudonyms.\todo{Check the four
descriptions against the attack definitions in \cite{kamel2020veremi} and the F2MD documentation; the wording below
follows our own inspection of the logs.}
\emph{GridSybil} (A16): the attacker broadcasts under several pseudonyms at the same time and claims positions a
few meters apart, which a receiver sees as a group of slow vehicles; each pseudonym is kept for many messages.
\emph{DataReplaySybil} (A17): the attacker re-broadcasts data taken from a genuine vehicle under a pseudonym that
it changes after one or two messages. \emph{DoSRandomSybil} (A18): the attacker sends at a higher rate than
normal, with random field values, and changes pseudonym with almost every message. \emph{DoSDisruptiveSybil}
(A19): the attacker sends at a higher rate, with content taken from other vehicles' messages, again under
pseudonyms used once or twice.

Two properties of this data matter for evaluation. First, a receiver can group messages only by pseudonym. For
A17 to A19 the sequence heard from one pseudonym is then one to three messages long in 95\% to 99.7\% of the
windows used in this paper, against 20\% for GridSybil and about 14\% for benign traffic. A short pseudonym life
is an observable signature of these attacks, and it is also a property that any sequence model can use without
modeling motion. Grouping by the true sender identifier of the simulation, as some studies do, removes this effect
and uses information that a receiver does not have. Second, neither dataset paper defines a train/test
split~\cite{heijden2018veremi,kamel2020veremi}, and the authors of VeReMi caution against using the data as the
sole basis for evaluating learned models~\cite{heijden2018veremi}. The successor dataset VeReMi NextGen ships
training, validation, and test sets~\cite{hermann2026veremi}. We further find that the attack scenarios of one
time window replay the same underlying traffic, so that a split made inside one scenario does not separate
vehicles across scenarios (Section~\ref{sec:protocol}).

The dataset paper reports a non-learned reference: plausibility and consistency checks combined by a fixed
threshold, applied to each received message~\cite{kamel2020veremi}. We re-implemented this detector offline and
obtain an F1 score within 0.05 of the published value on six of the eight Sybil subsets.\todo{Re-read Table~III of
\cite{kamel2020veremi}; the published values were taken from our notes.} Because it scores messages and the
detectors below score windows, it serves as context and is not part of the comparison.

\subsection{Learning-Based Detection on VeReMi}

Most learning-based detectors are supervised, as the survey by Boualouache and Engel
observes~\cite{boualouache2023survey}; it adds that the ability of the unsupervised ones to detect unseen attacks
is unverified, because they are evaluated on known attacks. Work on the original VeReMi addresses position
falsification with classical classifiers, often on plausibility
features~\cite{so2018integrating,sharma2021machine,ercan2022misbehavior,sharma2022machine,uprety2021privacy}.
On VeReMi-Extension, LSTM and CNN classifiers over message sequences~\cite{alladi2021securing,alladi2023deep},
a fine-tuned language model~\cite{hamhoum2024mistralbsm}, and recurrent networks on engineered
features~\cite{youness2025vemisnet,mahesh2025hyperlstm} report accuracy or F1 of 0.95 to 0.998 over many classes.
Lower values appear under more local or stricter settings: 0.68 to 0.80 when each vehicle learns from its own
messages~\cite{demir2025distributed}, and an F-score of 0.48 to 0.49 on DoS Sybil variants that were not in
training, 0.71 with knowledge transfer~\cite{sedar2024knowledge}.

Sybil attacks are the sole target of few studies. Kamel et al.\ classify the Sybil type from misbehavior reports
collected at a back-end authority~\cite{kamel2019misbehavior}; other work uses gradient boosting~\cite{fdil2024sybil}
or a graph network on the authors' own extended data~\cite{tang2025deep}. Elsewhere the Sybil attacks are some
classes among many, mostly the two DoS
variants~\cite{hamhoum2024mistralbsm,sedar2024knowledge,baharlouei2024evaluating}, and one study excludes them by
design~\cite{franceschini2025line}.

Split practice is summarized in Table~\ref{tab:related}. Where a split is described, it is drawn over windows or
sequences~\cite{youness2025vemisnet,franceschini2025line}, within each dataset or each vehicle's own
log~\cite{baharlouei2024evaluating,campos2025federated,demir2025distributed}, or by time and by
attack~\cite{sedar2024knowledge}; in many papers we could not find it. Among the papers we
read, we found none that reports a length or message-rate control, and none that discusses pseudonym lifetimes of
one or two messages.

\begin{table*}[t]
\caption{Closest prior work on VeReMi-Extension and this work. ``n.s.'': not stated in the text we could read;
``--'': not applicable.}
\label{tab:related}
\centering
\footnotesize
\setlength{\tabcolsep}{4pt}
\newcolumntype{L}[1]{>{\raggedright\arraybackslash}p{#1}}
\begin{tabular}{@{}L{2.4cm}L{2.1cm}L{3.5cm}L{2.9cm}L{2.8cm}L{2.2cm}@{}}
\toprule
Work & Learning setting & Unit & Split & Threshold source & Sybil evaluated \\
\midrule
Alladi et al.~\cite{alladi2021deep} & one-class (benign only) & 20 consecutive messages of one sender &
85\% of benign sequences for training; by vehicle n.s. & maximum accuracy on labeled test data & A16--A19 \\
PAMPOS~\cite{kalogiannis2026pampos} & one-class (benign only) & 10-message sliding window; sequences of
$\geq$\,15 messages; decision per sender & by sender, 70/15/15 & percentile of benign validation scores &
A16--A19 \\
MistralBSM~\cite{hamhoum2024mistralbsm} & supervised & windows of 10--100 messages grouped by sender ID & n.s. &
-- & A18, one of nine classes \\
VeMisNet~\cite{youness2025vemisnet} & supervised & message sequence, length n.s. & stratified 80/20 over
sequences & -- & among 20 classes; per attack n.s. \\
ADVENT eval.~\cite{baharlouei2024evaluating} & supervised & per-vehicle time-windowed statistics & 70/30 of each
dataset & n.s. & A18, A19 \\
Sedar et al.~\cite{sedar2024knowledge} & reinforcement learning, label-derived rewards & time series of
messages & by time and by attack & -- & A18, A19 (unseen) \\
\midrule
This work & benign only; unlabeled mix; supervised & receiver-side pseudonym link, up to 32 messages & by physical
vehicle, 80/10/10 & percentile of validation scores & A16--A19, each separately \\
\bottomrule
\end{tabular}
\end{table*}

\subsection{Reconstruction-Based and One-Class Detection}

One-class detectors are trained on benign sequences and flag inputs that they reconstruct or predict poorly.
Alladi et al.\ reconstruct sequences of 20 messages with CNN-LSTM and stacked LSTM networks and report an accuracy
of 0.98, with a recall of 0.98 to 1.0 on the four Sybil types; the threshold is the one with the highest accuracy
on the labeled test sequences~\cite{alladi2021deep}. The journal extension of this work reports the same
accuracy~\cite{alladi2021deepadv}. PAMPOS predicts the next message with a
causal transformer, splits by sender, takes the threshold from benign validation senders, and reports an F1 of 0.94
to 0.95 on A16 to A19; it requires sequences of at least 15 messages and decides per
sender~\cite{kalogiannis2026pampos}.\todo{Re-check the values of \cite{alladi2021deep} and
\cite{kalogiannis2026pampos} against the full texts before submission.} Further one-class or reconstruction models
use LSTM autoencoders~\cite{vishwanath2026federated}, generative adversarial
networks~\cite{zhong2023sybil,shahriar2024vehigan}, transformers~\cite{liu2023svmdformer}, and variational
autoencoders with Gaussian mixtures~\cite{campos2025federated}; in one case the reconstruction errors are passed
to an SVM that is fitted with labels~\cite{bao2026lstm}. In almost all of these, ``normal'' training data is
benign data selected with the ground-truth labels. Training on an unlabeled mix is rare: an LSTM autoencoder with a
Gaussian mixture model trained on mixed trajectory data reaches 51\% accuracy at its 99th-percentile threshold,
which the authors compare to random guessing~\cite{wang2024anomaly}.

Weak results are reported where an attack respects the physics of motion or replays valid data. In VehiGAN, every
model is at chance on a constant position offset, and a plain autoencoder on engineered features has a higher
average AUROC than the proposed ensemble~\cite{shahriar2024vehigan}; PAMPOS reports an F1 of 0.19 on the same
attack~\cite{kalogiannis2026pampos}. On the in-vehicle CAN bus, a full replay is detected at a rate of 0.57
because, as the authors note, replayed sequences are valid~\cite{longari2022candito}, and a plain window
autoencoder detects 3\% of replayed sequences~\cite{hanselmann2020canet}. For V2X Sybil attacks, we found no
comparable negative result and, to our knowledge, no study that compares several detectors under benign-only,
unlabeled-mix, and supervised training on one split.

\subsection{Warnings From Anomaly Detection and Security Research}

Several effects examined in this paper are known in general. Unsupervised anomaly detection assumes that the
training data is almost entirely normal; an autoencoder fits what is frequent in its training set, and methods for
contaminated data assume that anomalies are a small
minority~\cite{ruff2021unifying,pang2021deep,zhou2017anomaly,qiu2022latent}. Autoencoders and other generative
models can reconstruct anomalies well or score out-of-distribution inputs as more normal than the training
data~\cite{gong2019memorizing,yoon2021autoencoding,nalisnick2019deep,kirichenko2020normalizing,wang2019location}. Distances in a
latent space depend on clean, compact normal data~\cite{ruff2018deep,denouden2018improving}, and a small number of
labels can change results substantially~\cite{ruff2020deep,ruff2021rethinking,han2022adbench}. Models also exploit
shortcuts, features that separate the classes in a dataset without reflecting the intended
concept~\cite{geirhos2020shortcut,jacobs2022ai}.

On the evaluation side, time-series anomaly benchmarks have been criticized for trivial problems and unrealistic
anomaly density~\cite{wu2021current}, and simple baselines are often competitive with deep
models~\cite{kim2022rigorous,audibert2022deep,sarfraz2024position,schmidl2022anomaly}. Security research lists
sampling bias, data snooping, and inappropriate baselines among recurring
pitfalls~\cite{arp2022dos,sommer2010outside,pendlebury2019tesseract,apruzzese2023sok}; leakage between training
and test data is a documented source of irreproducible results~\cite{kapoor2023leakage}; and audits of intrusion
detection datasets show that dataset defects change reported
results~\cite{engelen2021troubleshooting,liu2022error,flood2024bad}. We do not present these effects as new. This
paper measures whether, and how strongly, they occur for Sybil detection on VeReMi-Extension.
